%-----------------------------------------------------------------------------------------
% Autor dieser Vorlage:
% Stefan Macke (http://fachinformatiker-anwendungsentwicklung.net)
% Permalink zur Vorlage: http://fiae.link/LaTeXVorlageFIAE
%
% Sämtliche verwendeten Abbildungen, Tabellen und Listings stammen von Dirk Grashorn.
%
% Lizenz: Creative Commons 4.0 Namensnennung - Weitergabe unter gleichen Bedingungen
% -----------------------------------------------------------------------------------------

\documentclass[
	ngerman,
	toc=listof, % Abbildungsverzeichnis sowie Tabellenverzeichnis in das Inhaltsverzeichnis aufnehmen
	toc=bibliography, % Literaturverzeichnis in das Inhaltsverzeichnis aufnehmen
	footnotes=multiple, % Trennen von direkt aufeinander folgenden Fußnoten
	parskip=half, % vertikalen Abstand zwischen Absätzen verwenden anstatt horizontale Einrückung von Folgeabsätzen
	numbers=noendperiod % Den letzten Punkt nach einer Nummerierung entfernen (nach DIN 5008)
]{scrartcl}
\pdfminorversion=5 % erlaubt das Einfügen von pdf-Dateien bis Version 1.7, ohne eine Fehlermeldung zu werfen (keine Garantie für fehlerfreies Einbetten!)
\usepackage[utf8]{inputenc}
\usepackage{glossaries} % muss als erstes eingebunden werden, da Meta/Packages ggfs. Sonderzeichen enthalten

% !TEX root = Projektdokumentation.tex

% Hinweis: der Titel muss zum Inhalt des Projekts passen und den zentralen Inhalt des Projekts deutlich herausstellen
\newcommand{\titel}{Workflow-Anpassungen}
\newcommand{\untertitel}{Anpassen von Workflows für das Änderungsmanagement in Windchill}
\newcommand{\kompletterTitel}{\titel{} -- \untertitel}

\newcommand{\autorName}{Fabian Schmidberger}
\newcommand{\autorAnschrift}{Am Galgenfeld 4a}
\newcommand{\autorOrt}{86554 Pöttmes}

\newcommand{\betriebLogo}{ArsandisLogo.pdf}
\newcommand{\betriebName}{Arsandis GmbH}
\newcommand{\betriebAnschrift}{Angkofen 5}
\newcommand{\betriebOrt}{85276 Pfaffenhofen an der Ilm}

\newcommand{\ausbildungsberuf}{Fachinformatiker/Fachinformatikerin (VO 2020) Fachrichtung: Anwendungsentwicklung}
\newcommand{\betreff}{Dokumentation zur betrieblichen Projektarbeit}
\newcommand{\pruefungstermin}{Winter 2023/2024}
\newcommand{\abgabeOrt}{Pfaffenhofen}
\newcommand{\abgabeTermin}{30.11.2023}
 % Metadaten zu diesem Dokument (Autor usw.)
% !TEX root = ../Projektdokumentation.tex

% Anpassung an Landessprache ---------------------------------------------------
\usepackage{babel}

% Mehr Linienoptionen für Tabellen
\usepackage{booktabs}

% Umlaute ----------------------------------------------------------------------
%   Umlaute/Sonderzeichen wie äüöß direkt im Quelltext verwenden (CodePage).
%   Erlaubt automatische Trennung von Worten mit Umlauten.
% ------------------------------------------------------------------------------
\usepackage[T1]{fontenc}
\usepackage{textcomp} % Euro-Zeichen etc.

% Schrift ----------------------------------------------------------------------
\usepackage{lmodern} % bessere Fonts
\usepackage{relsize} % Schriftgröße relativ festlegen

% Tabellen ---------------------------------------------------------------------
\PassOptionsToPackage{table}{xcolor}
\usepackage{tabularx}
% für lange Tabellen
\usepackage{longtable}
\usepackage{array}
\usepackage{ragged2e}
\usepackage{lscape}
\newcolumntype{w}[1]{>{\raggedleft\hspace{0pt}}p{#1}} % Spaltendefinition rechtsbündig mit definierter Breite

% Grafiken ---------------------------------------------------------------------
\usepackage[dvips,final]{graphicx} % Einbinden von JPG-Grafiken ermöglichen
\usepackage{graphics} % keepaspectratio
\usepackage{floatflt} % zum Umfließen von Bildern
\graphicspath{{Bilder/}} % hier liegen die Bilder des Dokuments

% Sonstiges --------------------------------------------------------------------
\usepackage[titles]{tocloft} % Inhaltsverzeichnis DIN 5008 gerecht einrücken
\usepackage{amsmath,amsfonts} % Befehle aus AMSTeX für mathematische Symbole
\usepackage{enumitem} % anpassbare Enumerates/Itemizes
\usepackage{xspace} % sorgt dafür, dass Leerzeichen hinter parameterlosen Makros nicht als Makroendezeichen interpretiert werden

\usepackage{makeidx} % für Index-Ausgabe mit \printindex
\usepackage[printonlyused]{acronym} % es werden nur benutzte Definitionen aufgelistet

% Einfache Definition der Zeilenabstände und Seitenränder etc.
\usepackage{setspace}
\usepackage{geometry}

% Symbolverzeichnis
\usepackage[intoc]{nomencl}
\let\abbrev\nomenclature
\renewcommand{\nomname}{Abkürzungsverzeichnis}
\setlength{\nomlabelwidth}{.25\hsize}
\renewcommand{\nomlabel}[1]{#1 \dotfill}
\setlength{\nomitemsep}{-\parsep}

\usepackage{varioref} % Elegantere Verweise. „auf der nächsten Seite“
\usepackage{url} % URL verlinken, lange URLs umbrechen etc.

\usepackage{chngcntr} % fortlaufendes Durchnummerieren der Fußnoten
% \usepackage[perpage]{footmisc} % Alternative: Nummerierung der Fußnoten auf jeder Seite neu

\usepackage{ifthen} % bei der Definition eigener Befehle benötigt
\usepackage{todonotes} % definiert u.a. die Befehle \todo und \listoftodos
\usepackage[square]{natbib} % wichtig für korrekte Zitierweise

% PDF-Optionen -----------------------------------------------------------------
\usepackage{pdfpages}
\pdfminorversion=5 % erlaubt das Einfügen von pdf-Dateien bis Version 1.7, ohne eine Fehlermeldung zu werfen (keine Garantie für fehlerfreies Einbetten!)
\usepackage[
    bookmarks,
    bookmarksnumbered,
    bookmarksopen=true,
    bookmarksopenlevel=1,
    colorlinks=true,
% diese Farbdefinitionen zeichnen Links im PDF farblich aus
    anchorcolor=AOBlau,% Ankertext
    citecolor=AOBlau, % Verweise auf Literaturverzeichniseinträge im Text
    filecolor=AOBlau, % Verknüpfungen, die lokale Dateien öffnen
    menucolor=AOBlau, % Acrobat-Menüpunkte
    urlcolor=AOBlau,
% diese Farbdefinitionen sollten für den Druck verwendet werden (alles schwarz)
    %linkcolor=black, % einfache interne Verknüpfungen
    %anchorcolor=black, % Ankertext
    %citecolor=black, % Verweise auf Literaturverzeichniseinträge im Text
    %filecolor=black, % Verknüpfungen, die lokale Dateien öffnen
    %menucolor=black, % Acrobat-Menüpunkte
    %urlcolor=black,
%
    %backref, % Quellen werden zurück auf ihre Zitate verlinkt
    pdftex,
    plainpages=false, % zur korrekten Erstellung der Bookmarks
    pdfpagelabels=true, % zur korrekten Erstellung der Bookmarks
    hypertexnames=false, % zur korrekten Erstellung der Bookmarks
    linkcolor=black,
    linktoc=all,
]{hyperref}
% Befehle, die Umlaute ausgeben, führen zu Fehlern, wenn sie hyperref als Optionen übergeben werden
\hypersetup{
    pdftitle={\titel -- \untertitel},
    pdfauthor={\autorName},
    pdfcreator={\autorName},
    pdfsubject={\titel -- \untertitel},
    pdfkeywords={\titel -- \untertitel},
}


% zum Einbinden von Programmcode -----------------------------------------------
\usepackage{listings}
\usepackage{xcolor}
\definecolor{hellgelb}{rgb}{1,1,0.9}
\definecolor{colKeys}{rgb}{0,0,1}
\definecolor{colIdentifier}{rgb}{0,0,0}
\definecolor{colComments}{rgb}{0,0.5,0}
\definecolor{colString}{rgb}{1,0,0}
\lstset{
    float=hbp,
	basicstyle=\footnotesize,
    identifierstyle=\color{colIdentifier},
    keywordstyle=\color{colKeys},
    stringstyle=\color{colString},
    commentstyle=\color{colComments},
    backgroundcolor=\color{hellgelb},
    columns=flexible,
    tabsize=2,
    frame=single,
    extendedchars=true,
    showspaces=false,
    showstringspaces=false,
    numbers=left,
    numberstyle=\tiny,
    breaklines=true,
    breakautoindent=true,
	captionpos=b,
}
\lstdefinelanguage{cs}{
	sensitive=false,
	morecomment=[l]{//},
	morecomment=[s]{/*}{*/},
	morestring=[b]",
	morekeywords={
		abstract,event,new,struct,as,explicit,null,switch
		base,extern,object,this,bool,false,operator,throw,
		break,finally,out,true,byte,fixed,override,try,
		case,float,params,typeof,catch,for,private,uint,
		char,foreach,protected,ulong,checked,goto,public,unchecked,
		class,if,readonly,unsafe,const,implicit,ref,ushort,
		continue,in,return,using,decimal,int,sbyte,virtual,
		default,interface,sealed,volatile,delegate,internal,short,void,
		do,is,sizeof,while,double,lock,stackalloc,
		else,long,static,enum,namespace,string},
}
\lstdefinelanguage{natural}{
	sensitive=false,
	morecomment=[l]{/*},
	morestring=[b]",
	morestring=[b]',
	alsodigit={-,*},
	morekeywords={
		DEFINE,DATA,LOCAL,END-DEFINE,WRITE,CALLNAT,PARAMETER,USING,
		IF,NOT,END-IF,ON,*ERROR-NR,ERROR,END-ERROR,ESCAPE,ROUTINE,
		PERFORM,SUBROUTINE,END-SUBROUTINE,CONST,END-FOR,END,FOR,RESIZE,
		ARRAY,TO,BY,VALUE,RESET,COMPRESS,INTO,EQ},
}
\lstdefinelanguage{php}{
	sensitive=false,
	morecomment=[l]{/*},
	morestring=[b]",
	morestring=[b]',
	alsodigit={-,*},
	morekeywords={
		abstract,and,array,as,break,case,catch,cfunction,class,clone,const,
		continue,declare,default,do,else,elseif,enddeclare,endfor,endforeach,
		endif,endswitch,endwhile,extends,final,for,foreach,function,global,
		goto,if,implements,interface,instanceof,namespace,new,old_function,or,
		private,protected,public,static,switch,throw,try,use,var,while,xor
		die,echo,empty,exit,eval,include,include_once,isset,list,require,
		require_once,return,print,unset},
}
 % verwendete Packages
% !TEX root = ../Projektdokumentation.tex

% Seitenränder -----------------------------------------------------------------
\setlength{\topskip}{\ht\strutbox} % behebt Warnung von geometry
\geometry{a4paper,left=25mm,right=20mm,top=20mm,bottom=20mm}

\usepackage[
	automark, % Kapitelangaben in Kopfzeile automatisch erstellen
	headsepline, % Trennlinie unter Kopfzeile
	ilines % Trennlinie linksbündig ausrichten
]{scrlayer-scrpage}

% Kopf- und Fußzeilen ----------------------------------------------------------
\pagestyle{scrheadings}
% chapterpagestyle gibt es nicht in scrartcl
%\renewcommand{\chapterpagestyle}{scrheadings}
\clearscrheadfoot

% Kopfzeile
\renewcommand{\headfont}{\normalfont} % Schriftform der Kopfzeile
\ihead{\large{\textsc{\titel}}\\ \small{\untertitel} \\[2ex] \textit{\headmark}}
\chead{}
\ohead{\includegraphics[scale=0.2]{\betriebLogo}}
\setlength{\headheight}{15mm} % Höhe der Kopfzeile
%\setheadwidth[0pt]{textwithmarginpar} % Kopfzeile über den Text hinaus verbreitern (falls Logo den Text überdeckt)

% Fußzeile
\ifoot{\autorName}
\cfoot{}
\ofoot{\pagemark}

% Überschriften nach DIN 5008 in einer Fluchtlinie
% ------------------------------------------------------------------------------

% Abstand zwischen Nummerierung und Überschrift definieren
% > Schön wäre hier die dynamische Berechnung des Abstandes in Abhängigkeit
% > der Verschachtelungstiefe des Inhaltsverzeichnisses
\newcommand{\headingSpace}{1.5cm}

% Abschnittsüberschriften im selben Stil wie beim Inhaltsverzeichnis einrücken
\renewcommand*{\othersectionlevelsformat}[3]{
  \makebox[\headingSpace][l]{#3\autodot}
}

% Für die Einrückung wird das Paket tocloft benötigt
%\cftsetindents{chapter}{0.0cm}{\headingSpace}
\cftsetindents{section}{0.0cm}{\headingSpace}
\cftsetindents{subsection}{0.0cm}{\headingSpace}
\cftsetindents{subsubsection}{0.0cm}{\headingSpace}
\cftsetindents{figure}{0.0cm}{\headingSpace}
\cftsetindents{table}{0.0cm}{\headingSpace}


% Allgemeines
% ------------------------------------------------------------------------------

\onehalfspacing % Zeilenabstand 1,5 Zeilen
\frenchspacing % erzeugt ein wenig mehr Platz hinter einem Punkt

% Schusterjungen und Hurenkinder vermeiden
\clubpenalty = 10000
\widowpenalty = 10000
\displaywidowpenalty = 10000

% Quellcode-Ausgabe formatieren
\lstset{numbers=left, numberstyle=\tiny, numbersep=5pt, breaklines=true}
\lstset{emph={square}, emphstyle=\color{red}, emph={[2]root,base}, emphstyle={[2]\color{blue}}}

\counterwithout{footnote}{section} % Fußnoten fortlaufend durchnummerieren
\setcounter{tocdepth}{3} % im Inhaltsverzeichnis werden die Kapitel bis zum Level der subsubsection übernommen
\setcounter{secnumdepth}{3} % Kapitel bis zum Level der subsubsection werden nummeriert

% Aufzählungen anpassen
\renewcommand{\labelenumi}{\arabic{enumi}.}
\renewcommand{\labelenumii}{\arabic{enumi}.\arabic{enumii}.}
\renewcommand{\labelenumiii}{\arabic{enumi}.\arabic{enumii}.\arabic{enumiii}}

% Tabellenfärbung:
\definecolor{heading}{rgb}{0.64,0.78,0.86}
\definecolor{odd}{rgb}{0.9,0.9,0.9}
 % Definitionen zum Aussehen der Seiten
\input{Allgemein/Befehle} % eigene allgemeine Befehle, die z.B. die Arbeit mit LaTeX erleichtern
\input{Befehle} % eigene projektspezifische Befehle, z.B. Abkürzungen usw.

\begin{document}


%\phantomsection
%\thispagestyle{empty}
%\pdfbookmark[1]{Eidesstattliche Erklärung}{ihkdeckblatt}
%\includegraphicsKeepAspectRatio{DeckblattIHK}{1}
%\cleardoublepage

\phantomsection
\thispagestyle{plain}
\pdfbookmark[1]{Deckblatt}{deckblatt}
% !TEX root = Projektdokumentation.tex
\begin{titlepage}

\begin{center}
\includegraphics[scale=0.25]{LogoIHK.pdf}\\[1ex]
\Large{Abschlussprüfung \pruefungstermin}\\[3ex]

\Large{\ausbildungsberuf}\\
\LARGE{\betreff}\\[4ex]

\huge{\textbf{\titel}}\\[1.5ex]
\Large{\textbf{\untertitel}}\\[4ex]

\normalsize
Abgabetermin: \abgabeOrt, den \abgabeTermin\\[3em]
\textbf{Prüfungsbewerber:}\\
\autorName\\
\autorAnschrift\\
\autorOrt\\
Prüflingsnummer: 20525\\[5ex]

\includegraphics[scale=0.4]{\betriebLogo}\\[2ex]
\textbf{Ausbildungsbetrieb:}\\
\betriebName\\
\betriebAnschrift\\
\betriebOrt\\
Ansprechpartner: Joachim Loos\\[5em]
\end{center}

\end{titlepage}
\cleardoublepage

% Preface --------------------------------------------------------------------
\phantomsection
\pagenumbering{Roman}
\pdfbookmark[1]{Inhaltsverzeichnis}{inhalt}
\tableofcontents

\cleardoublepage

\phantomsection
\listoffigures
\cleardoublepage

\phantomsection
\listoftables
\cleardoublepage

\phantomsection
\lstlistoflistings
\cleardoublepage

\newcommand{\abkvz}{Abkürzungsverzeichnis}
\renewcommand{\nomname}{\abkvz}
\section*{\abkvz}
\markboth{\abkvz}{\abkvz}
\addcontentsline{toc}{section}{\abkvz}
% !TEX root = Projektdokumentation.tex

% Es werden nur die Abkürzungen aufgelistet, die mit \ac definiert und auch benutzt wurden. 
%
% \acro{VERSIS}{Versicherungsinformationssystem\acroextra{ (Bestandsführungssystem)}}
% Ergibt in der Liste: VERSIS Versicherungsinformationssystem (Bestandsführungssystem)
% Im Text aber: \ac{VERSIS} -> Versicherungsinformationssystem (VERSIS)

% Hinweis: allgemein bekannte Abkürzungen wie z.B. bzw. u.a. müssen nicht ins Abkürzungsverzeichnis aufgenommen werden
% Hinweis: allgemein bekannte IT-Begriffe wie Datenbank oder Programmiersprache müssen nicht erläutert werden,
%          aber ggfs. Fachbegriffe aus der Domäne des Prüflings (z.B. Versicherung)

% Die Option (in den eckigen Klammern) enthält das längste Label oder
% einen Platzhalter der die Breite der linken Spalte bestimmt.
\begin{acronym}[WWWWW]
	\acro{CSV}{Comma Separated Value}
	\acro{PLM}{Product Lifecycle Management}
	\acro{AWS}{Applied Water Systems}
	\acro{UML}{Unified Modeling Language}
	\acro{RDP}{Remote Desktop Protocol}
	\acro{XGV}{Xylem Global Vault}
	\acro{ECR}{Engineering Change Request}
	\acro{ECN}{Engineering Change Notice}
	\acro{ECM}{Engineering Change Management}
	\acro{ECA}{Engineering Change Activity}
	\acro{JSP}{Jakarta Server Pages}
	\acro{OOTB}{out of the box}
	\acro{PTC}{Parametric Technology Corporation}
	\acro{PR}{Promotion Request}

\end{acronym}

\clearpage

% Inhalt ---------------------------------------------------------------------
\pagenumbering{arabic}
% !TEX root = Projektdokumentation.tex
% !TEX root = ../Projektdokumentation.tex
\section{Einleitung}
\label{sec:Einleitung}

\subsection{Vorstellung der Arsandis GmbH}
\label{sec:Vorstellung_Arsandis}
Die Arsandis GmbH ist ein IT- und Dienstleistungsunternehmen in Angkofen, Pfaffenhofen an der Ilm.
Gegründet wurde die Arsandis im Jahr 2015 und beschäftigt zur Zeit 13 Mitarbeiter.

Das Hauptgeschäftsfeld der Arsandis GmbH liegt in der Beratung von Fertigungsunternehmen im Bereich des Product Lifecycle Managements (PLM).
Darüber hinaus hat das Unternehmen ein starkes Interesse an zukunftsorientierten Technologien und verfolgt aktiv die Entwicklungen in Bereichen wie Augmented Reality, Virtual Reality und Internet of Things.
Diese Technologien bieten innovative Lösungen für die Industrie der Zukunft.
Weitere zukunftsorientierte Technologien werden ebenfalls verfolgt, darunter zum Beispiel AR-, VR- und IoT-basierte Lösungen für die Industrie der Zukunft.
Die Arsandis GmbH hat ein breites Spektrum an Kunden aus verschiedenen Branchen.
Dazu gehören Unternehmen aus dem Maschinenbau, der Automobil- und Luftfahrtindustrie sowie der HighTech- und Telekommunikationsbranche.

Als offizieller Partner von PTC liegt das Hauptaugenmerk hier meist auf dem PTC-Portfolio, das ein umfangreiches Ökosystem für die Fertigungsindustrie bereitstellt.

\subsection{Projektumfeld}
\label{sec:Projektumfeld}
Das Projekt wird im Auftrag der amerikanischen Firma \ac{AWS} durchgeführt, die in der Herstellung von Wasserpumpen tätig ist.
\ac{AWS} ist ein Tochterunternehmen von Xylem, einem Unternehmen, das sich auf die Entwicklung und Bereitstellung von Wassertechnologien spezialisiert hat.
\ac{AWS} nutzt den Windchill Server von Xylem, der als \ac{XGV} bekannt ist.
Die meisten Tochterunternehmen von Xylem, einschließlich \ac{AWS}, sind auf diesem Server als Organisation eingerichtet.

Durchgeführt wird das Projekt von mir in den Räumen der Arsandis GmbH.
Im Vornherein wurde für dieses Projekt eine Windchill Entwicklungsumgebung auf unserem Unternehmensserver eingerichtet.
Auf dieser findet jegliche Entwicklungsarbeit statt.
Nach Abschluss der Implementierungsphase sollen die Softwareanpassungen in das Testsystem von Xylem integriert werden.
Die finalen Tests und das Deployment auf das Produktivsystem werden von Mitarbeitern von AWS durchgeführt.

\subsection{Projektziel} 
\label{sec:Projektziel}
Das Projekt verfolgt zwei Ziele.
Zum einen soll das bestehende \ac{ECM} von \ac{AWS} angepasst und erweitert werden.
Dem \ac{ECM}-Prozess sollen neue Funktionen hinzugefügt werden, die die Robustheit des Prozesses erhöhen und die benötigte Durchlaufzeit verringern.

Andererseits sollen neue Organisationen in das Windchill-\ac{ECM} eingebunden werden.
Dafür soll der manuelle Prozess der Befüllung von Excel-Dateien durch Windchill ersetzt werden.
User sollen den \ac{ECM}-Prozess in Windchill initiieren, wodurch alle Daten, die für das \ac{ECM} relevant sind, gebündelt an einem Ort abgelegt werden.
Da sich die CAD-Dateien ohnehin schon auf dem Windchill-Server befinden, lassen sich CAD-Dateien und \ac{ECM} somit praktisch verknüpfen.

\subsection{Projektbegründung} 
\label{sec:Projektbegruendung}
Es wird eine erhöhte Robustheit des aktuellen Prozesses angestrebt, indem das manuelle Erstellen von \acp{PR} automatisiert wird, was auch zu einer Zeitersparnis führt.
Außerdem soll der \ac{ECM}-Prozess in Organisationen eingeführt werden, die bisher entweder auf Excel-Tabellen und selbst entwickelte Webanwendungen oder auf persönliche Besprechung mit dem Produktmanager zurückgegriffen haben, um diese Prozesse abzubilden.
Dies führt zum einen zu einer erheblichen Zeitersparnis, da viele manuelle Vorgänge, wie das Befüllen von einer Excel-Tabelle, wegfallen.
Dadurch wird wiederum auch die Robustheit deutlich erhöht, da repetitive Tätigkeiten nicht mehr vergessen werden können und häufige Fehler vermieden werden können.

\subsection{Projektabgrenzung}
Das Abschlussprojekt nimmt ungefähr 80 \% des kompletten Projekts ein.
Die restlichen 20 \% werden nach Abschluss des Abschlussprojekts angegangen.
Dann ist geplant, noch einige weitere Quality-Of-Life Verbesserungen zu implementieren.

\subsection{Projektschnittstellen}
\label{sec:Projektschnittstellen}

\subsubsection{Windchill}
Windchill ist eine \ac{PLM} Webanwendung von PTC, die das Verwalten von Unternehmensobjekten für die Herstellungsindustrie vereinfacht.
Besonders Wert wird dabei auf die Verwaltung von CAD-Modellen gelegt.
CAD-Modelle sind am Computer gefertigte zwei- und dreidimensionale Zeichnungen von Konstruktionsobjekten.
Mit dem Windchill Workgroup Manager existiert eine Schnittstelle zwischen Windchill und vielen CAD-Anwendungen, die auf dem Markt existieren (wie \zB Creo Parametric, CATIA oder Solidworks).

\subsubsection{\acs{XGV}}
Der \acl{XGV} ist der Windchill-Server von Xylem.
Die meisten Tochterfirmen von Xylem sind auf diesem Server als eigenständige Organisation eingerichtet.
Für die Durchführung des Projekts hat dieser Server allerdings keine Relevanz, da die endgültige Übernahme auf das Produktivsystem nicht von uns durchgeführt wird.
Dennoch wurde er aus Gründen der Vollständigkeit aufgeführt.

\subsubsection{\ac{XGV}-Trainingsserver}
Dieses Testsystem ist ein Klon vom \ac{XGV}-Produktivserver.
Hier nimmt der Kunde die betriebseigenen Tests vor, um unsere Anpassungen zu überprüfen.
Zusammen mit meinem Projektleiter werde ich die Übernahme von unserem Entwicklungssystem auf den Trainingsserver durchführen.

\subsubsection{\acs{AWS}}
Auch \ac{AWS} hat eigene Tochterunternehmen, die in Windchill als Organisation eingerichtet sind.
Darunter sind zum Beispiel FHD, Vogel oder Lowara.
Wichtig ist, dass alle Änderungen, die von uns vorgenommen werden, ausschließlich die Organisationen von \ac{AWS} betreffen.

\subsubsection{Workflows}
Workflows sind Objekte in Windchill, mit denen man Unternehmensprozesse abbilden kann.
Dafür stellt Windchill einen Workflow-Editor bereit, mit dem man die Unternehmensprozesse in einer praktischen Benutzeroberfläche definieren kann.
Innerhalb von Workflows lässt sich auch direkt Code definieren, der dann mit dem Workflow ausgeführt wird.
Hier kann man auch auf Klassen und Methoden innerhalb der Windchill-Codebase verweisen.
Die folgenden Konzepte wurden und werden mit Hilfe von Workflows umgesetzt.

\subsubsection{\acl{ECM}}
Das Änderungsmanagement ist ein essenzieller Bestandteil von Windchill.
Teile, wie \zB CAD-Dateien oder Dokumente kommen in ihrem Lebenszyklus zumeist irgendwann an den Punkt, an dem sie für die Produktion freigegeben werden.
Da bereits umfassende Ressourcen für die Produktion aufgewendet wurden, muss eine nachträgliche Änderung gut begründet werden.
Schließlich wurden bereits Maschinen/Maschinenteile bestellt, eingerichtet oder modifiziert, um das Produkt herzustellen.
Demzufolge müssen nachträgliche Änderungen nun in Absprache mit verschiedenen Unternehmensabteilungen koordiniert werden.
Der Change Admin ist für die Koordination des Change Managements verantwortlich.
In Windchill gibt es eigentlich drei verschiedene Change Administrator Rollen, vereinfacht wird folgend aber lediglich vom Change Admin gesprochen.

Hier kommt das Änderungsmanagement ins Spiel.
Es umfasst insgesamt drei Bestandteile, die alle in Windchill implementiert sind: \acl{ECR}, \acl{ECN} und \acl{ECA}.
Diese werden in den nächsten Abschnitten genauer erläutert.

\subsubsection{\acl{ECR}}
Der Änderungsantrag ist der erste Schritt im Änderungsmanagement.
Grob gesagt wird im \ac{ECR} geklärt, \underline{ob} die Änderung umgesetzt werden kann.
Der Change Admin wählt die Unternehmensabteilungen aus, die über die Durchführung der Änderung abstimmen sollen.
Die Abteilungen stimmen dann entweder für oder gegen die geplante Änderung.
Nur wenn alle befragten Abteilungen für die Änderung sind, wird der Prozess mit der \ac{ECN} fortgesetzt.
Stimmt mindestens eine Abteilung dagegen, so wird der Änderungsantrag verworfen.

\subsubsection{\acl{ECN}}
Die Änderungsmitteilung ist der zweite Schritt im Änderungsmanagement.
Hier geht es darum, dass festgelegt wird, \underline{wie} die Änderung umgesetzt wird.
Dafür legt der Change-Admin zuerst fest, welche Abteilungen an der Umsetzung der Änderungen beteiligt sein sollen.
Danach bestimmt er, welche Nutzer aus der Abteilung an der Umsetzung der Änderung arbeiten sollen.

\subsubsection{\acl{ECA}}
Die Änderungsaufgabe ist der dritte und letzte Schritt im Änderungsmanagement.
Die geplanten Änderungen werden hier umgesetzt und anschließend vom Change-Admin überprüft.
Ist der Change-Admin unzufrieden mit den umgesetzten Änderungen, so kann er eine Überarbeitung von den relevanten Abteilungen anfordern.
Sobald der Change-Administrator seine Zustimmung zu den Änderungen erteilt hat, markiert dies den erfolgreichen Abschluss der Änderungsaufgabe und damit auch des gesamten Change-Prozesses.

\subsubsection{\acl{PR}}
Der Erhöhungsantrag ist nicht direkt Teil dieses Projekts.
Allerdings werden neue Mechanismen im \ac{ECM}-Workflow eingeführt, die das manuelle Erstellen von Erhöhungsanträgen ersetzen soll.
Demzufolge wird der Erhöhungsantrag kurz erläutert.

Ein Windchill-Objekt durchläuft stets verschiedene Phasen seines Lebenszyklus.
Im Moment der Erstellung befindet es sich üblicherweise im Status `In Arbeit`.
Um das Objekt in einen anderen Zustand zu überführen, wie beispielsweise `Freigegeben`, wird ein Erhöhungsantrag durchlaufen.
In diesem Prozess entscheidet der Genehmiger, ob eine Änderung des Status des Objekts angebracht ist.
% !TEX root = ../Projektdokumentation.tex
\section{Projektplanung} 
\label{sec:Projektplanung}


\subsection{Projektphasen}
\label{sec:Projektphasen}

Für die Umsetzung des Projekts wurden 76 Stunden angesetzt.
Der Start des Projekts erfolgt am 07.11.2023 und bis zum 27.11.2023 wird es abgeschlossen sein.
Der Tabelle \ref{tab:Zeitplanung} kann entnommen werden, in welche Hauptphasen das Projekt gegliedert wurde.
\tabelle{Zeitplanung}{tab:Zeitplanung}{ZeitplanungKurz}\\
Eine detailliertere Zeitplanung findet sich im \Anhang{app:Zeitplanung}.
Im Vergleich zum Projektantrag gab es einige Änderungen.
Zum einen wurde im Antrag vergessen die Erstellung der Projektdokumentation zu listen.
Dadurch erhöht sich die Gesamtdauer des Projekts auf 80 Stunden und die für die Erstellung eines Benutzerhandbuchs geschätzte Zeit wurde auf 1 Stunde reduziert.


\subsection{Ressourcenplanung}
\label{sec:Ressourcenplanung}

Für die Durchführung des Projekts werden folgende Ressourcen verwendet.

\subsubsection{Hardware}
Von meinem Büroarbeitsplatz aus wird mein Windows 10 Arbeitslaptop genutzt.

\subsubsection{Entwicklungsumgebung}
Im Vorfeld des Projekts wurde bereits eine virtuelle Maschine mit dem Betriebssystem Windows Server 2019 aufgesetzt.
Diese ist mit 16 GB RAM ausgestattet.
Auf dieser wurde Windchill in der Version 11.1 M020-CPS026 installiert.
Windchill benötigt eine Oracle Datenbank der Version 19C.
Windchill benötigt eine Java Runtime Environment der Version 1.8.0 Update 202.
Für die Erstellung des Codes wird die IDE Eclipse auf der \acs{VM} verwendet.
Die User des Windchillsystems werden im Open Source LDAP Programm OpenDJ 3.0.0 verwaltet.
Eine weitere Komponente, die benötigt wird, ist ein Apache 2.4 Webserver.

\subsubsection{Testsystem}
Das Testsystem läuft auf einem Linux-Server ohne User Interface.
Für die Verbindung zur Weboberfläche und zum Linux-Server wird das VPN Programm GlobalProtect genutzt.

\subsubsection{Weitere Software}
Zur Erstellung der Projektdokumentation wird IntelliJ mit dem TeXiFy Plugin und die MiKTeX Distribution für LaTeX verwendet.

\subsubsection{Personal}
Außerdem haben mich folgende Personen bei meinem Projekt unterstützt.
Der Product Owner von \acs{AWS} legt die Anforderungen fest und nimmt das Projekt ab.
Ein Experte für das Windchill-System von \acs{AWS} hilft mir dabei, mich in die Customizations von \acs{AWS} einzuarbeiten.
Als Entwickler führt der Autor die Umsetzung des Projektes durch.
Der Projektleiter von Arsandis überprüft die Umsetzung und den Code.
Eine Anwendungsspezialistin von Arsandis unterstützt den Autor beim Implementieren und Testen.
Eine Mitarbeitern der Personalabteilung hilft beim Aufstellen der Wirtschaftlichkeitsanalyse.

\subsection{Entwicklungsprozess}
\label{sec:Entwicklungsprozess}

Die Umsetzung meines Projekts wird grundsätzlich agil ablaufen, in dem der aktuelle Status und Änderungen an den Anforderungen nach Bedarf mit den Kunden diskutiert werden.
Außerdem gibt es regelmäßige Besprechungen mit dem Projektleiter, um den Status des Projekts zu besprechen, dort wird auch diskutiert, ob es noch Optimierungsbedarf bei der Implementierung gibt.

% !TEX root = ../Projektdokumentation.tex
\section{Analysephase} 
\label{sec:Analysephase}


\subsection{Ist-Analyse} 
\label{sec:IstAnalyse}
Dieser Abschnitt erhält Auszüge aus dem Projektantrag, um eine Übersicht über die aktuelle Situation des Systems zu geben.

Der \ac{ECM} besteht aus den zwei Komponenten \ac{ECR} und \ac{ECN}.
Aktuell werden Aufgaben im \ac{ECN} getrackt, dies ist so von Windchill eigentlich nicht vorgesehen.
In Zukunft sollen Aufgaben deswegen in einem \ac{ECA} abgebildet werden.
Dies führt dazu, dass die konkreten Implementierungsaufgaben logisch von der Planungsphase (ECN) abgegrenzt wird, was es dem Change Admin einfacher macht, einen Überblick über die aktuelle Situation zu bekommen.

Im \ac{ECM} von \ac{AWS} sollen neue Abteilungen berücksichtigt werden.
Dafür müssen diese dem bestehenden \ac{ECR} und \ac{ECN} hinzugefügt werden.

Bei der Erstellung der \acsp{ECR} werden zudem Startparameter abgefragt.
Das sind Boolean-Attribute, die die automatische Auswahl von Abteilungen beeinflussen.
Dies ist eine Erleichterung für den Change Admin, da dann nicht jedes Mal alle Abteilungen manuell ausgefüllt werden müssen.
Der \ac{ECR} ist für alle Organisationen gleich.
Allerdings benötigen nicht alle Organisationen alle Attribute, die auf dem \ac{ECR} definiert wurden.
Dafür wurde von einem anderen Softwareunternehmen eine XML-Datei entwickelt, die nicht benötigte Attribute von einer Organisation vom UI versteckt.
Es sollen neue Boolean-Attribute definiert werden, die nur für die \ac{AWS} Organisationen sichtbar sind.
Dazu müssen diese dem \ac{ECR} hinzugefügt werden und die in der XML-Datei für die anderen Organisationen versteckt werden.

Um den Change Admin bei der Auswahl der beteiligten Abteilungen zu unterstützen, wurde von einem anderen Softwareunternehmen eine benutzerdefinierte JSP Seite entwickelt, die die Abteilungen in einer Liste darstellt.
Der Change Admin kann dann über Checkboxen Abteilungen für den Prozess entfernen oder hinzufügen.
Dieses Pop-Up-Fenster wurde von einem anderen Softwareunternehmen in einem bereits abgeschlossenem Projekt entwickelt.


Zudem werden Abteilungen aktuell in einem eigens dafür angefertigten Pop-Up-Fenster ausgewählt.
Hier kann der Change Admin über Checkboxen Abteilungen auswählen, die am Change Prozess beteiligt sein sollen.
Dieses Pop-Up-Fenster wurde von einem anderen Softwareunternehmen in einem bereits abgeschlossenem Projekt entwickelt.
Allerdings bringt diese Lösung einige Probleme mit sich.
Zum einen verzögern Customizations, die von den OOTB-Windchill-Konfigurationen abweichen, oft größere Windchill Upgrades, da nicht garantiert ist, dass die Customization auch in zukünftigen Versionen funktioniert.
Zum anderen können nur Administratoren mit Zugriff zur virtuellen Maschine Änderungen an der Struktur des Pop-Up-Fensters vornehmen (\zB hinzufügen oder entfernen von Abteilungen), da die Customization als JSP-Datei auf dem Server liegt.
Um diese beiden Probleme zu beseitigen, soll hier auf eine Lösung zurückgegriffen werden, die bereits von Haus aus in Windchill implementiert ist, nämlich den 'Set Up Participants'-Tab.
Da es sich um eine OOTB-Funktion von Windchill handelt, ist es unwahrscheinlich, dass es damit Probleme bei einem Upgrade gibt.
Außerdem können Windchill-Administratoren über den Workflow Template Editor Änderungen an den Abteilungen vornehmen, ohne Zugriff auf die VM zu benötigen.

Zudem sind die aktuellen Workflowprozesse sehr unstrukturiert dargestellt und sind kaum dokumentiert.
Dies macht es schwieriger für Entwickler, die nicht mit den Prozessen vertraut sind, Erweiterungen oder Änderungen vorzunehmen.
\subsection{Wirtschaftlichkeitsanalyse}
\label{sec:Wirtschaftlichkeitsanalyse}

Da die exakten Gesamtkosten des Projekts nicht offenbart wurden, gehen wir von einem empfohlenen Festpreis von 50.000 € aus, wie mir mitgeteilt wurde. %TODO: Silvia fragen, wegen Projektkosten
Der Anteil dieser spezifischen Projektarbeit am Gesamtprojekt wird auf etwa 30 \% geschätzt.
Daher kalkulieren wir einen Umsatz von 15.000 € für die Umsetzung gemäß der hier beschriebenen Aufgabenstellung.

\subsubsection{Projektkosten}
\subsubsection*{Ausbildungsvergütung}
Hier haben wir ein Bruttoentgelt von 1080 €.

\subsubsection*{Arbeitgeberanteil Sozialversicherung}
Die Arbeitgeber-Sozialabgaben betragen 199,92 € ohne die Unfallversicherung.
%todo: Silvia fragen

\subsubsection*{Abstimmungskosten}
Der Projektleiter von Arsandis führte mehrere interne Abstimmungsgespräche zum Fortschritt des Projekts durch.
Diese Gespräche hatten das Ziel sicherzustellen, dass die Qualitätsanforderungen erfüllt werden und die Umsetzung auf bewährten Methoden basiert, während die Lösung sich gut in den Gesamtentwurf des Kundensystems einfügt.
Zusätzlich nahm der Projektleiter, wie in anderen Teilen des Projekts auch, an Meetings mit dem Kunden teil. Konkrete Zahlen sind für mich nicht einsehbar.
Es wurde mir geraten, mit einem festen Stundensatz von 100 € zu kalkulieren.

\subsubsection*{Gemeinkosten}
Intern verwenden wir keinen Gemeinkostenfaktor.
Ungeachtet dessen müssen wir jedoch sicherstellen, dass unsere Gemeinkosten gedeckt werden.
Aus diesem Grund werde ich eine weitere Annahme treffen, um eine standardmäßige Wirtschaftlichkeitsanalyse durchzuführen, da mir die genauen Zahlen nicht zugänglich sind.
\begin{itemize}
	\item Marketingkosten: Abgedeckt durch eine Mitarbeiterin = 15 \%
	\item Energiekosten: Heizung, Strom = 15 \%
	\item Mietkosten: Büroräume = 15 \%
	\item Ausrüstungskosten: Clients, Monitore, Server, Peripherie, Lizenzen = 15 \%
\end{itemize}

Das führt uns zu einem Gemeinkostenzuschlag von 60 \%

\subsubsection*{Gesamtkosten}
Hier werden die Kosten eines Arbeitstages durch das Herunterbrechen des Jahresverdiensts auf einen Tag ermittelt.
\Zitat{ArbeitstageImJahr}{b}Für das Jahr 2023 wurden 251 Arbeitstage angenommen.

Kosten pro Arbeitstag = 12960 € / (251 Arbeitstage - 24 Urlaubstage) = 57,09 €



\subsubsection{Amortisation}

\subsection{Anwendungsfälle}
\label{sec:Anwendungsfaelle}
\begin{itemize}
	\item Welche Anwendungsfälle soll das Projekt abdecken?
	\item Einer oder mehrere interessante (!) Anwendungsfälle könnten exemplarisch durch ein Aktivitätsdiagramm oder eine EPK detailliert beschrieben werden.
\end{itemize}

\paragraph{Beispiel}
Ein Beispiel für ein Use Case-Diagramm findet sich im \Anhang{app:UseCase}.


\subsection{Qualitätsanforderungen}
\label{sec:Qualitaetsanforderungen}
\begin{itemize}
	\item Welche Qualitätsanforderungen werden an die Anwendung gestellt (\zB hinsichtlich Performance, Usability, Effizienz \etc )?
\end{itemize}


\subsection{Lastenheft/Fachkonzept}
\label{sec:Lastenheft}
Zu Beginn des Projekts hat der Kunde bereits einen Entwurf des Lastenhefts ausgearbeitet, welcher zusammen mit dem Autor und dem Projektleiter von Arsandis weiter angepasst wurde.
Dort sind alle Anforderungen des Auftraggebers für die Anpassungen der aktuellen Anwendung enthalten.
Im \Anhang{app:Lastenheft} befindet sich ein Auszug aus dem Lastenheft.

% !TEX root = ../Projektdokumentation.tex
\section{Entwurfsphase} 
\label{sec:Entwurfsphase}

\subsection{Zielplattform}
\label{sec:Zielplattform}

\begin{itemize}
	\item Beschreibung der Kriterien zur Auswahl der Zielplattform (\ua Programmiersprache, Datenbank, Client/Server, Hardware).
\end{itemize}

\subsection{Geschäftslogik}
\label{sec:Geschaeftslogik}

\begin{itemize}
	\item Modellierung und Beschreibung der wichtigsten (!) Bereiche der Geschäftslogik (\zB mit Kom\-po\-nen\-ten-, Klassen-, Sequenz-, Datenflussdiagramm, Programmablaufplan, Struktogramm, EPK).
	\item Wie wird die erstellte Anwendung in den Arbeitsfluss des Unternehmens integriert?
\end{itemize}


\subsection{Maßnahmen zur Qualitätssicherung}
\label{sec:Qualitaetssicherung}
\begin{itemize}
	\item Welche Maßnahmen werden ergriffen, um die Qualität des Projektergebnisses (siehe Kapitel~\ref{sec:Qualitaetsanforderungen}: \nameref{sec:Qualitaetsanforderungen}) zu sichern (\zB automatische Tests, Anwendertests)?
	\item \Ggfs Definition von Testfällen und deren Durchführung (durch Programme/Benutzer).
\end{itemize}

Das Testing erfolgt manuell durch \glqq durchklicken\grqq{} der Workflowprozesse und wird in zwei Phasen aufgeteilt.
Zuerst erfolgt ein Arsandis-Interner Test auf unserer Entwicklungsumgebung, bei dem eine Kollegin den Prüfling unterstützt.
Sobald wir unseren Testprozess erfolgreich abgeschlossen haben, werden die Anpassungen auf das Testsystem des Kunden übertragen.
Daran anschließend startet die zweite Testphase, die von einigen Mitarbeitern von \ac{AWS} durchgeführt werden.
Nach ihren Tests erhalten wir ein Testprotokoll, in dem alle Fehler beziehungsweise Auffälligkeiten dokumentiert wurden.

Ein großer Teil des Codes wird direkt in den Workflows geschrieben, wodurch sich der Einsatz von Unit-Tests als sehr schwierig bis unmöglich gestaltet, weswegen es keine automatisierten Tests gibt.

\subsection{Pflichtenheft/Datenverarbeitungskonzept}
\label{sec:Pflichtenheft}
\begin{itemize}
	\item Auszüge aus dem Pflichtenheft/Datenverarbeitungskonzept, wenn es im Rahmen des Projekts erstellt wurde.
\end{itemize}

\paragraph{Beispiel}
Ein Beispiel für das auf dem Lastenheft (siehe Kapitel~\ref{sec:Lastenheft}: \nameref{sec:Lastenheft}) aufbauende Pflichtenheft ist im \Anhang{app:Pflichtenheft} zu finden.

% !TEX root = ../Projektdokumentation.tex
\section{Implementierungsphase}  %todo: finish section
\label{sec:Implementierungsphase}

\subsubsection{Anpassen der Workflow-Templates}
\subsubsection*{Umstrukturierung der Nodes zur besseren Übersichtlichkeit}
Exemplarisch wird hier nur die neue Struktur des ECRs gezeigt, da sich hier die größten Änderungen ergeben haben.
Der ECN Workflow hat allerdings auch eine Umstrukturierung erfahren.
Es wurde besonders Wert darauf gelegt, dass die Elemente in einer klaren Struktur angeordnet werden und sich Linien möglichst nicht kreuzen.
Man kann deutlich erkennen, dass der Workflow nun deutlich aufgeräumter und strukturierter wirkt.
Dadurch wird vor allem das Leseverständnis erhöht.
Zudem wurde der große Block im unteren Teil von \autoref{WorkflowVergleich1} mit einem seperaten Block (Department Approval Process) ersetzt.

Ein vorher-nachher Vergleich findet sich im \Anhang{WorkflowVergleich}.



\subsubsection{Ergänzung von Lebenszyklusstatus im ECA}
\label{LebenszyklusstatusECA}

\subsubsection{Erstellen neuer Types}
Damit der Workflow für ECA-E und ECA-NE auch im System getriggert wird, müssen 

\begin{itemize}
	\item Anlegen neuer Types
	\subitem OIR
	\subitem Lifecycle Template
	\item Anpassen der Lifecycle Templates
	\item Erstellung der Java-Klassen
\end{itemize}


% !TEX root = ../Projektdokumentation.tex
\section{Abnahmephase} %todo: finish section
\label{sec:Abnahmephase}

\begin{itemize}
	\item Welche Tests (\zB Unit-, Integrations-, Systemtests) wurden durchgeführt und welche Ergebnisse haben sie geliefert (\zB Logs von Unit Tests, Testprotokolle der Anwender)?
	\item Wurde die Anwendung offiziell abgenommen?
\end{itemize}

\paragraph{Beispiel}
Ein Auszug eines Unit Tests befindet sich im \Anhang{app:Test}. Dort ist auch der Aufruf des Tests auf der Konsole des Webservers zu sehen.

% !TEX root = ../Projektdokumentation.tex
\section{Einführungsphase}
\label{sec:Einfuehrungsphase}

\begin{itemize}
	\item Welche Schritte waren zum Deployment der Anwendung nötig und wie wurden sie durchgeführt (automatisiert/manuell)?
	\item Wurden \ggfs Altdaten migriert und wenn ja, wie?
	\item Wurden Benutzerschulungen durchgeführt und wenn ja, Wie wurden sie vorbereitet?
\end{itemize}

% !TEX root = ../Projektdokumentation.tex
\section{Dokumentation}
\label{sec:Dokumentation}

\begin{itemize}
	\item Wie wurde die Anwendung für die Benutzer/Administratoren/Entwickler dokumentiert (\zB Benutzerhandbuch, API-Dokumentation)?
	\item Hinweis: Je nach Zielgruppe gelten bestimmte Anforderungen für die Dokumentation (\zB keine IT-Fachbegriffe in einer Anwenderdokumentation verwenden, aber auf jeden Fall in einer Dokumentation für den IT-Bereich).
\end{itemize}

\paragraph{Beispiel}
Ein Ausschnitt aus der erstellten Benutzerdokumentation befindet sich im \Anhang{app:BenutzerDoku}.
Die Entwicklerdokumentation wurde mittels PHPDoc\footnote{Vgl. \cite{phpDoc}} automatisch generiert.
Ein beispielhafter Auszug aus der Dokumentation einer Klasse findet sich im \Anhang{app:Doc}.

% !TEX root = ../Projektdokumentation.tex
\section{Fazit} 
\label{sec:Fazit}

\subsection{Soll-/Ist-Vergleich}
\label{sec:SollIstVergleich}

Das Projektziel wurde zur Deadline nicht eingehalten.
Grund dafür ist, dass sich auf der einen Seite die Anforderungen im Vergleich zum Start geändert haben.
Auf der anderen Seite war die Zeit von rund 80 Stunden deutlich zu knapp bemessen für die Fertigstellung des Projekts.
Darum wurde die Deadline um einige Wochen nach hinten verschoben.
Das hat auch zur Folge, dass das Projektbudget aufgestockt werden musste.

Durch den Zeitdruck und eine suboptimale Kommunikation sind außerdem einige Fehler in der Implementierung passiert.
Teilweise wurden Anforderungen nicht optimal umgesetzt oder anders umgesetzt, als sich der Kunde das vorgestellt hat.
Die Fehler wurden von den Testern entdeckt und dokumentiert.
Das Projekt wird auch über das Abschlussprojekt hinaus forgesetzt, dabei sollen die Fehler ausgebessert werden.
Trotzdem ist der Auftraggeber grundsätzlich zufrieden mit dem Ergebnis, denn es soll noch ein weiteres Projekt folgen.

Auch das programmatische Umbenennen von Rollen wurde in der geplanten Zeit nicht geschafft.
Deswegen wurde diese Aufgabe an einen anderen Mitarbeiter der Arsandis ausgelagert.
Dies ist aber nicht weiter tragisch, da diese Funktion erst für den Umzug auf den Produktivserver wirklich relevant ist.
Für das Testen auf den Testservern wurden einige Team-Templates händisch angepasst, wodurch diese Phase des Projekts nicht beeinflusst wurde.

Im Projektantrag wurde bei den Aufgaben die Erstellung eines Use-Case Diagramms erwähnt, dies wurde während des Projekts allerdings fallen gelassen, da es für das Projekt nicht notwendig erschien.

Abschließend ist der Zeitaufwand für das Anfertigen der Dokumentation unerwartet hoch ausgefallen, dafür wurde weniger Zeit in das Erstellen der Aktivitätsdiagramme investiert, damit 80 Stunden nicht überschritten werden.

Wie in Tabelle~\ref{tab:Vergleich} zu erkennen ist, konnte die Zeitplanung nicht wie erwartet eingehalten werden.
\tabelle{Soll-/Ist-Vergleich}{tab:Vergleich}{Zeitnachher.tex}


\subsection{Lessons Learned}
\label{sec:LessonsLearned}

Dieses Projekt war eine hervorragende Möglichkeit, das bis dahin in der Ausbildung Gelernte anzuwenden.
Da dies das erste große Projekt war, in dem ich mich mit der Zeitplanung hauptsächlich selbst befassen sollte, war dies zu Beginn eine große Herausforderung.
Allerdings waren die Schulstunden zum Thema Projektmanagement eine große Hilfe, um Kundenanforderungen richtig zu planen.

Zudem habe ich durch den regelmäßigen Austausch mit dem Kunden viel im Bereich der Kommunikation dazugewonnen.
Es wurden zwar bereits vor dem Abschlussprojekt regelmäßige Gespräche mit langjährigen Kunden geführt, die Kunden dieses Projekts kannten wir allerdings noch nicht zuvor.
Teilweise gab es Missverständnisse, die nach einem klärenden Gespräch beseitigt wurden.
Dies war eine wertvolle Erfahrung, die im zukünftigen Berufsleben Gold wert ist.


% Literatur ------------------------------------------------------------------
\clearpage
\renewcommand{\refname}{Literaturverzeichnis}
\bibliography{Bibliographie}
\bibliographystyle{Allgemein/natdin} % DIN-Stil des Literaturverzeichnisses
% !TEX root = Projektdokumentation.tex
\clearpage
\addsec{Eidesstattliche Erklärung}

% Hinweis: die eidesstattliche Erklärung ist ggfs. an die Richtlinie der IHK anzupassen

Ich, \autorName, versichere hiermit, dass ich meine \textbf{\betreff} mit dem
Thema
\begin{quote}
\textit{\kompletterTitel}
\end{quote}
selbständig verfasst und keine anderen als die angegebenen Quellen und Hilfsmittel benutzt habe,
wobei ich alle wörtlichen und sinngemäßen Zitate als solche gekennzeichnet habe. Die Arbeit
wurde bisher keiner anderen Prüfungsbehörde vorgelegt und auch nicht veröffentlicht.\\[6ex]

\abgabeOrt, den \abgabeTermin

\includegraphics[width=0.5\textwidth]{Unterschrift.png}

\textsc{\autorName}


% Anhang ---------------------------------------------------------------------
\clearpage
\appendix
\pagenumbering{roman}
% !TEX root = Projektdokumentation.tex
\section{Anhang}
\subsection{Detaillierte Zeitplanung}
\label{app:Zeitplanung}

\tabelleAnhang{ZeitplanungKomplett}


\end{document}
